
\chapter*{Microexperimento IC10: cambiar realmente el tokenizador}
\phantomsection\addcontentsline{toc}{chapter}{Microexperimento IC10: tokenizadores comparables}
\label{lab:ic10}
\begin{idea}[Ahora si cambiamos la representacion, no el estimador de probabilidades]
En el laboratorio anterior se compararon \emph{dos bigramas de caracteres}
con distinta procedencia de datos. Aquello ilustra estimacion de
probabilidades. Aqui comparamos dos \emph{tokenizadores BPE} que aprenden
fusiones diferentes, uno en prosa general y otro en prosa tecnica, con el
\textbf{mismo presupuesto de 20 fusiones}. Esta es una primera medicion de
la variable de representacion que estudia IC10, sin afirmar que un
modelo de lenguaje mejora por ello.
\end{idea}
\section*{Reglas transparentes}
La unidad inicial de cada palabra es un punto de código Unicode; se añade
un marcador \texttt{</w>} para no fusionar entre palabras. Cada paso combina
el par de unidades adyacentes mas frecuente del corpus de entrenamiento
(con desempate lexicografico). Las reglas se aplican, en el orden en que
se aprendieron, a las \textbf{mismas seis frases tecnicas de prueba}.
La separacion entrenamiento/prueba es por documento. Se conserva la
normalizacion explicita: minusculas, recorte externo de espacios y
separacion por blancos. No se eliminan acentos, ni se usa el texto de
prueba para elegir fusiones. Este BPE didactico no es el tokenizador de
un modelo neuronal preentrenado.
\section*{Resultado verificado en Python, datos sinteticos}
\begin{center}\begin{tabular}{lrrr}\toprule
Tokenizador & Palabras de prueba & Tokens en prueba & Tokens/palabra\\\midrule
BPE entrenado con prosa general & 35 & 148 & 4.228571\\
BPE entrenado con texto tecnico & 35 & 127 & 3.628571\\\bottomrule
\end{tabular}\end{center}
La diferencia agregada es $127-148=-21$ tokens y
$3.628571-4.228571=-0.600000$ tokens por palabra, bajo estas reglas.
La tabla por documento se genera en
\nolinkurl{experimentos/resultados_ic10_tokenizacion.csv}; las veinte
fusiones de cada modelo, en \nolinkurl{experimentos/fusiones_ic10.csv}.
\begin{verificacion}[Interpretacion responsable]
El resultado no permite afirmar que un modelo general tenga peor
perplejidad ni que el tokenizador especializado mejore una tarea aplicada.
Para contrastar \textbf{bits por caracter} entre tokenizadores se
necesita asignar probabilidades normalizadas y compatibles a las mismas
secuencias (incluyendo convenciones de fronteras y fin), y evaluar en
textos reservados. Menos tokens significa \emph{otra segmentacion}; no
implica por si solo menor entropia cruzada ni mejor generalizacion.
\end{verificacion}
\section*{Ejecucion y reporte}
\begin{verbatim}
python codigo/ic10_tokenizadores.py
# En un equipo con R instalado:
Rscript codigo/ic10_tokenizadores.R
python verificaciones/comparar_r_python.py
\end{verbatim}
El estudiante reporta las dos comparaciones por separado:
\textbf{A)} corpus de entrenamiento frente a BPC de bigramas por
caracteres; \textbf{B)} entrenamiento del tokenizador frente a tokens por
palabra. Ambas son ilustrativas y no equivalen a un estudio industrial.
